\chapter{\MakeUppercase{Systems of Nonlinear Equations}}

\section{Multivariable Newton Method}

Newton's method can be extended to systems of nonlinear equations in a natural way. Consider a system of the form
\begin{equation}
F(\mathbf{x}) = \mathbf{0},
\end{equation}
where \( F : \mathbb{R}^n \to \mathbb{R}^n \) is a vector-valued function.

Starting from an initial guess \( \mathbf{x}_0 \), the method generates a sequence \( \{\mathbf{x}_k\} \) using the iterative formula
\begin{equation}
\mathbf{x}_{k+1} = \mathbf{x}_k + \mathbf{s}_k,
\end{equation}
where the step \( \mathbf{s}_k \) is obtained by solving the linear system
\begin{equation}
J_F(\mathbf{x}_k)\mathbf{s}_k = -F(\mathbf{x}_k).
\end{equation}

Here, \( J_F(\mathbf{x}_k) \) denotes the Jacobian matrix evaluated at \( \mathbf{x}_k \). This formulation can be interpreted as a multivariable generalization of the Newton update, where the nonlinear system is locally approximated by a linear one.

At each iteration, a linear system must be solved. Therefore, compared to the scalar case, the computational cost per iteration is higher. However, under suitable conditions, the method retains its rapid convergence properties.

If the initial guess is sufficiently close to the solution and the Jacobian matrix is nonsingular at the root, the method typically exhibits quadratic convergence. This makes it one of the most efficient methods for solving nonlinear systems \cite{kelley, burden}.

Despite its strong convergence properties, the method has several limitations. The need to compute and factorize the Jacobian matrix can be computationally expensive, especially for large systems. In addition, if the Jacobian is singular or nearly singular, the method may fail or produce unstable steps.

For these reasons, modifications and alternative methods are often considered in practice, particularly when dealing with large-scale or ill-conditioned systems.

\paragraph{Implementation Example.}

The following implementation shows the core computational structure of Newton's method for nonlinear systems. At each iteration, the Jacobian system is solved in order to compute the Newton step.

\begin{lstlisting}[caption={Python implementation of Newton's method for systems}, label={lst:newton_system}]
def newton_system(F, J, x0, tol=1e-8, max_iter=100):

    x = np.array(x0, dtype=float)
    history = []

    for i in range(max_iter):

        Fx = F(x)
        norm_Fx = np.linalg.norm(Fx)

        if norm_Fx <= tol:
            return x, history

        Jx = J(x)
        s = np.linalg.solve(Jx, -Fx)

        x_next = x + s
        step_norm = np.linalg.norm(s)

        history.append({
            "iteration": i,
            "x": x.tolist(),
            "norm_f": norm_Fx,
            "step_norm": step_norm
        })

        if step_norm <= tol:
            return x_next, history

        x = x_next

    return x, history
\end{lstlisting}

\section{Convergence Analysis}

The convergence behavior of Newton's method for systems is similar to the scalar case, but it depends on additional factors related to the structure of the system.

Let \( \mathbf{x}^* \) be a solution of the system \( F(\mathbf{x}) = \mathbf{0} \). If the function \( F \) is sufficiently smooth and the Jacobian matrix \( J_F(\mathbf{x}^*) \) is nonsingular, then Newton's method exhibits quadratic convergence in a neighborhood of \( \mathbf{x}^* \). This means that the error satisfies
\begin{equation}
\|\mathbf{x}_{k+1} - \mathbf{x}^*\| \approx C \|\mathbf{x}_k - \mathbf{x}^*\|^2,
\end{equation}
for some constant \( C > 0 \) and sufficiently large \( k \) \cite{kelley}.

This rapid convergence makes Newton's method very efficient when a good initial approximation is available. However, the method is not globally convergent. If the initial guess is far from the solution, the method may fail to converge or may converge to an unintended solution.

Another important factor is the conditioning of the Jacobian matrix. If \( J_F(\mathbf{x}_k) \) is ill-conditioned or nearly singular, solving the linear system
\begin{equation}
J_F(\mathbf{x}_k)\mathbf{s}_k = -F(\mathbf{x}_k)
\end{equation}
may lead to large or unstable steps. This can cause divergence or oscillatory behavior.

In practice, the convergence of the method is often monitored using a combination of criteria such as
\begin{equation}
\|\mathbf{F}(\mathbf{x}_k)\| < \varepsilon
\end{equation}
and
\begin{equation}
\|\mathbf{x}_{k+1} - \mathbf{x}_k\| < \varepsilon,
\end{equation}
where \( \varepsilon \) is a prescribed tolerance.

Due to these limitations, various modifications of Newton's method have been developed to improve its robustness. These include damping strategies, line search techniques, and quasi-Newton methods such as Broyden's method.

\section{Broyden's Method}

Broyden's method is a quasi-Newton method developed for solving systems of nonlinear equations. The main idea of the method is to avoid the explicit computation of the Jacobian matrix while still retaining a structure similar to Newton's method.

In Newton's method, the step \( \mathbf{s}_k \) is obtained by solving
\begin{equation}
J_F(\mathbf{x}_k)\mathbf{s}_k = -F(\mathbf{x}_k).
\end{equation}
In contrast, Broyden's method replaces the Jacobian matrix with an approximation \( B_k \), leading to the system
\begin{equation}
B_k \mathbf{s}_k = -F(\mathbf{x}_k).
\end{equation}

The approximation \( B_k \) is updated at each iteration using information from previous steps. One common update formula, known as the good Broyden update, is given by
\begin{equation}
B_{k+1} = B_k + \frac{(\mathbf{y}_k - B_k \mathbf{s}_k)\mathbf{s}_k^T}{\mathbf{s}_k^T \mathbf{s}_k},
\end{equation}
where
\begin{equation}
\mathbf{y}_k = F(\mathbf{x}_{k+1}) - F(\mathbf{x}_k).
\end{equation}

This update ensures that the new approximation \( B_{k+1} \) satisfies a secant-like condition, which is a generalization of the secant method to higher dimensions.

One of the main advantages of Broyden's method is that it eliminates the need to compute the Jacobian matrix explicitly. This can significantly reduce computational cost, especially for large systems or when evaluating derivatives is expensive.

However, the convergence of Broyden's method is generally slower than that of Newton's method. While Newton's method achieves quadratic convergence under suitable conditions, Broyden's method typically exhibits superlinear convergence.

The performance of the method also depends on the choice of the initial approximation \( B_0 \). A good initial approximation can improve convergence speed, while a poor choice may lead to slow convergence or instability.

Despite these limitations, Broyden's method is widely used in practice due to its efficiency and its ability to handle problems where the Jacobian is difficult to compute \cite{kelley}.

\paragraph{Implementation Example.}

The following implementation fragment illustrates the main structure of Broyden's method. Instead of recomputing the exact Jacobian matrix at every iteration, the method updates an approximation using a rank-one correction formula.

\begin{lstlisting}[caption={Python implementation of Broyden's method}, label={lst:broyden}]
def broyden(F, x0, B0=None, tol=1e-8, max_iter=100):

    x = np.array(x0, dtype=float)

    if B0 is None:
        B = np.eye(len(x))
    else:
        B = np.array(B0, dtype=float)

    Fx = F(x)
    history = []

    for i in range(max_iter):

        s = np.linalg.solve(B, -Fx)

        x_next = x + s
        Fx_next = F(x_next)

        history.append({
            "iteration": i,
            "x": x.tolist(),
            "norm_f": np.linalg.norm(Fx),
            "step_norm": np.linalg.norm(s)
        })

        if np.linalg.norm(Fx_next) <= tol:
            return x_next, history

        y = Fx_next - Fx

        B += np.outer((y - B @ s), s) / (s @ s)

        x = x_next
        Fx = Fx_next

    return x, history
\end{lstlisting}

\section{Illustrative Example}

To illustrate the behavior of the methods for systems of nonlinear equations, consider the following system:
\begin{equation}
\begin{cases}
x^2 + y^2 - 4 = 0, \\
x - y - 1 = 0.
\end{cases}
\end{equation}

This system represents the intersection of a circle and a line. One of its solutions is approximately
\[
\mathbf{x}^* \approx (1.82288,\ 0.82288).
\]

The same system is used to demonstrate both Newton's method and Broyden's method. The goal is to observe how the iterates approach the solution.

\subsection{Newton Method for Systems}

Consider the nonlinear system
\begin{equation}
\begin{cases}
x^2 + y^2 - 4 = 0, \\
x - y - 1 = 0.
\end{cases}
\end{equation}

The initial approximation is selected as
\[
\mathbf{x}_0=(2,1).
\]

This starting point is chosen because it lies relatively close to the expected intersection region of the circle and the line. In Newton-type methods, selecting an initial approximation near the solution generally improves convergence behavior and reduces the risk of divergence.

The vector-valued function is
\[
F(x,y)=
\begin{bmatrix}
x^2+y^2-4 \\
x-y-1
\end{bmatrix},
\]
and the Jacobian matrix is
\[
J_F(x,y)=
\begin{bmatrix}
2x & 2y \\
1 & -1
\end{bmatrix}.
\]

At the initial point
\[
(2,1),
\]
the function value becomes
\[
F(2,1)=
\begin{bmatrix}
2^2+1^2-4 \\
2-1-1
\end{bmatrix}
=
\begin{bmatrix}
1 \\
0
\end{bmatrix}.
\]

The Jacobian matrix evaluated at the initial point is
\[
J_F(2,1)=
\begin{bmatrix}
4 & 2 \\
1 & -1
\end{bmatrix}.
\]

Newton's method requires solving the linear system
\[
J_F(\mathbf{x}_0)\mathbf{s}_0
=
-F(\mathbf{x}_0).
\]

Thus,
\[
\begin{bmatrix}
4 & 2 \\
1 & -1
\end{bmatrix}
\begin{bmatrix}
s_1 \\
s_2
\end{bmatrix}
=
\begin{bmatrix}
-1 \\
0
\end{bmatrix}.
\]

From the second equation,
\[
s_1-s_2=0,
\]
so
\[
s_1=s_2.
\]

Substituting into the first equation:
\[
4s_1+2s_1=-1.
\]

Therefore,
\[
6s_1=-1,
\qquad
s_1=s_2=-\frac16.
\]

The Newton step is therefore
\[
\mathbf{s}_0=
\left(
-\frac16,
-\frac16
\right).
\]

Updating the approximation:
\[
\mathbf{x}_1
=
\mathbf{x}_0+\mathbf{s}_0
=
\left(
2-\frac16,
1-\frac16
\right).
\]

Thus,
\[
\mathbf{x}_1
=
(1.83333,\ 0.83333).
\]

Evaluating the residual norm:
\[
\|F(\mathbf{x}_1)\|
\approx
0.05556.
\]

At the second iteration, the Jacobian matrix becomes
\[
J_F(1.83333,0.83333)
=
\begin{bmatrix}
3.66666 & 1.66666 \\
1 & -1
\end{bmatrix}.
\]

Solving the corresponding Newton system produces the next approximation
\[
\mathbf{x}_2
=
(1.82292,\ 0.82292),
\]
with residual norm
\[
\|F(\mathbf{x}_2)\|
\approx
0.00022.
\]

A third iteration gives
\[
\mathbf{x}_3
=
(1.82288,\ 0.82288),
\]
where the residual becomes extremely small.

Table 4.1 summarizes the iteration sequence.

\begin{table}[h!]
\centering
\caption{Newton method for the nonlinear system}
\label{tab:newton_system_example}
\begin{tabular}{|c|c|c|c|}
\hline
Iteration & \(x_k\) & \(y_k\) & \(\|F(\mathbf{x}_k)\|\) \\ \hline
0 & 2.00000 & 1.00000 & 1.00000 \\ \hline
1 & 1.83333 & 0.83333 & 0.05556 \\ \hline
2 & 1.82292 & 0.82292 & 0.00022 \\ \hline
3 & 1.82288 & 0.82288 & 0.00000 \\ \hline
\end{tabular}
\end{table}

The numerical results demonstrate the rapid convergence behavior of Newton's method for nonlinear systems. Once the iteration approaches the solution, the residual norm decreases very quickly, which is consistent with the quadratic convergence properties discussed earlier.

\subsection{Broyden's Method for Systems}

For Broyden's method, the same nonlinear system is considered:
\begin{equation}
\begin{cases}
x^2 + y^2 - 4 = 0, \\
x - y - 1 = 0.
\end{cases}
\end{equation}

The initial approximation is again selected as
\[
\mathbf{x}_0=(2,1).
\]

As in Newton's method, this point is chosen because it lies relatively close to the expected solution and produces a well-defined Jacobian structure.

Unlike Newton's method, Broyden's method does not recompute the exact Jacobian matrix at every iteration. Instead, the method constructs an approximation that is updated iteratively.

For the initial approximation of the Jacobian matrix, the exact Jacobian evaluated at the initial point is used:
\[
B_0=
\begin{bmatrix}
4 & 2 \\
1 & -1
\end{bmatrix}.
\]

This choice provides a stable starting approximation and improves the early convergence behavior of the method.

At the initial iteration, the nonlinear system is
\[
F(2,1)=
\begin{bmatrix}
1 \\
0
\end{bmatrix}.
\]

The step is obtained by solving
\[
B_0\mathbf{s}_0
=
-F(\mathbf{x}_0).
\]

Thus,
\[
\begin{bmatrix}
4 & 2 \\
1 & -1
\end{bmatrix}
\begin{bmatrix}
s_1 \\
s_2
\end{bmatrix}
=
\begin{bmatrix}
-1 \\
0
\end{bmatrix}.
\]

This produces the same first step as Newton's method:
\[
\mathbf{s}_0=
\left(
-\frac16,
-\frac16
\right).
\]

Therefore,
\[
\mathbf{x}_1
=
(1.83333,\ 0.83333).
\]

The residual norm becomes
\[
\|F(\mathbf{x}_1)\|
\approx
0.05556.
\]

At this stage, instead of recomputing the exact Jacobian, Broyden's method updates the approximation matrix using the vectors
\[
\mathbf{s}_0
=
\mathbf{x}_1-\mathbf{x}_0
\]
and
\[
\mathbf{y}_0
=
F(\mathbf{x}_1)-F(\mathbf{x}_0).
\]

The update formula is
\[
B_{1}
=
B_0
+
\frac{
(\mathbf{y}_0-B_0\mathbf{s}_0)\mathbf{s}_0^T
}{
\mathbf{s}_0^T\mathbf{s}_0
}.
\]

Using the updated approximation matrix, the next step is computed and the iteration proceeds.

At the second iteration, the approximation becomes
\[
\mathbf{x}_2
=
(1.82353,\ 0.82353),
\]
with residual norm
\[
\|F(\mathbf{x}_2)\|
\approx
0.00346.
\]

A third iteration gives
\[
\mathbf{x}_3
=
(1.82288,\ 0.82288),
\]
where the residual norm becomes very small.

Table 4.2 summarizes the iterations.

\begin{table}[h!]
\centering
\caption{Broyden method for the nonlinear system}
\label{tab:broyden_example}
\begin{tabular}{|c|c|c|c|}
\hline
Iteration & \(x_k\) & \(y_k\) & \(\|F(\mathbf{x}_k)\|\) \\ \hline
0 & 2.00000 & 1.00000 & 1.00000 \\ \hline
1 & 1.83333 & 0.83333 & 0.05556 \\ \hline
2 & 1.82353 & 0.82353 & 0.00346 \\ \hline
3 & 1.82288 & 0.82288 & 0.00001 \\ \hline
4 & 1.82288 & 0.82288 & 0.00000 \\ \hline
\end{tabular}
\end{table}

Compared to Newton's method, Broyden's method converges more slowly because it uses approximate Jacobian information. However, the method avoids repeated Jacobian evaluations, which may significantly reduce computational cost in large-scale problems or applications where derivative calculations are expensive.

\subsection{Comparison for the Example}

The results show that Newton's method converges faster than Broyden's method in this example. This is expected, since Newton's method uses exact derivative information, while Broyden's method relies on approximations.

On the other hand, Broyden's method avoids the computation of the Jacobian matrix, which can be advantageous in problems where derivative evaluation is expensive.

\subsection{Failure Cases}

\paragraph{Newton Method.}
Newton's method may fail when the Jacobian matrix becomes singular or nearly singular. Consider the system
\begin{equation}
\begin{cases}
x^2 = 0, \\
y = 0.
\end{cases}
\end{equation}

The Jacobian matrix is
\[
J_F(x,y) =
\begin{bmatrix}
2x & 0 \\
0 & 1
\end{bmatrix}.
\]

At the point \( (0, y) \), the Jacobian is singular. Starting from
\[
\mathbf{x}_0 = (0, 2),
\]
we have
\[
F(\mathbf{x}_0) = (0, 2),
\qquad
J_F(\mathbf{x}_0) =
\begin{bmatrix}
0 & 0 \\
0 & 1
\end{bmatrix}.
\]

The Newton step requires solving
\[
J_F(\mathbf{x}_0)\mathbf{s}_0 = -F(\mathbf{x}_0),
\]
which becomes
\[
\begin{bmatrix}
0 & 0 \\
0 & 1
\end{bmatrix}
\mathbf{s}_0 =
\begin{bmatrix}
0 \\
-2
\end{bmatrix}.
\]

This system has no unique solution, since the Jacobian matrix is singular. Therefore, the method cannot proceed from this initial point.

Even when the Jacobian is not exactly singular but close to singular, the computed step may become very large, leading to instability and divergence.

\paragraph{Broyden Method.}
Broyden's method may fail when the update of the Jacobian approximation becomes numerically unstable. Consider again the system
\begin{equation}
\begin{cases}
x^2 + y^2 - 4 = 0, \\
x - y - 1 = 0,
\end{cases}
\end{equation}
and assume that two consecutive iterates are very close:
\[
\mathbf{x}_k = (1.82288,\ 0.82288), \qquad
\mathbf{x}_{k+1} = (1.82289,\ 0.82289).
\]

Then the step
\[
\mathbf{s}_k = \mathbf{x}_{k+1} - \mathbf{x}_k
\]
is extremely small. In the Broyden update formula
\[
B_{k+1} = B_k + \frac{(\mathbf{y}_k - B_k \mathbf{s}_k)\mathbf{s}_k^T}{\mathbf{s}_k^T \mathbf{s}_k},
\]
the denominator \( \mathbf{s}_k^T \mathbf{s}_k \) becomes very small.

As a result, the update term may become very large, causing a significant distortion in the Jacobian approximation. This can lead to unstable iterations and loss of convergence.

In practice, such situations are handled by introducing safeguards, such as skipping the update or restarting the method with a new approximation.

\section{Comparison of Methods}

There are trade-offs between these two methods.

Newton's method is usually more accurate and converges more quickly provided that the Jacobian matrix is well-conditioned and the starting point is close enough to the root. The requirement to calculate the Jacobian and solve systems using the Jacobian matrix is computationally expensive.

In Broyden's method, on the contrary, the computational overhead is avoided by not calculating the Jacobian explicitly.

In general, Newton's method is better than Broyden's when higher accuracy and speed of convergence are essential. Conversely, Broyden's method is preferred when computational efficiency is crucial \cite{kelley, chapra}.